
% ============================================================
% File: 1_FRONTMATTER/05_preface.tex
% Book: Ontology of Viability
% ============================================================

\cleardoublepage

\chapter*{Preface}
\addcontentsline{toc}{chapter}{Preface}

\section*{The Place of This Book}

This book is the first foundational volume of the
\textit{Philosophy of Viability}. Its purpose is to establish
the ontological foundation upon which the broader philosophical
system can be developed.

The project begins with a fundamental question: what does it
mean for something to be? Before viability can be examined as
a condition, process, capacity, or mode of existence, the
ontological conditions of Being itself must be investigated.

Any coherent account of viability presupposes an account of
what exists, how entities possess identity, how they become
determinate, how their structures and relations are constituted,
and how they manifest themselves within Being. It also
presupposes an examination of the conditions under which
entities undergo change, maintain continuity, and enter into
different forms of ontological determination.

These questions are not preliminary matters that can be
disregarded once a definition of viability has been introduced.
They constitute the philosophical ground upon which the
concept of viability must be established.

Accordingly, this volume does not claim to present a complete
theory of viability. Its task is more fundamental: to construct
the ontological framework within which viability can be
defined, examined, and progressively developed.

\section*{The Ontological Starting Point}

The central starting point of this volume is the inseparable
relation between Being and identity.

Identity is not treated as an additional property assigned
to an entity after its existence has already been constituted.
An entity does not first exist as an ontologically neutral
something and subsequently acquire identity as an external
determination.

Nor are Being and identity regarded as two independently
constituted realities that become related only through a
subsequent connection.

Rather, Being and identity are ontologically inseparable,
while remaining analytically distinguishable. This distinction
allows their respective dimensions to be examined without
turning their conceptual differentiation into an ontological
separation.

Neither Being nor identity is granted ontological priority
over the other. Their relation is constitutive of the
ontological reality under investigation, rather than an
accidental relation between two independently complete
entities.

This commitment does not mean that every question concerning
identity has already been resolved. On the contrary, it
establishes the problem that guides the first part of this
book: how is identity related to determination, unity,
distinction, individuation, numerical identity, continuity,
and manifestation?

These concepts must be examined in their respective
ontological functions. Identity cannot simply be equated
with any one of them, nor can their relations be assumed
without argument.

The ontological constitution of identity, its relation to
the different realms of Being, and the dynamics through
which it is maintained or transformed require a broader
investigation. The first part therefore establishes the
conceptual distinctions necessary for the subsequent
development of the ontology.

\section*{The Independence of Being from the Knower}

A fundamental methodological commitment of this volume
is that Being is investigated independently of the
knowing subject.

The existence of an entity is not constituted by its
being perceived, recognized, described, or understood.
Its identity is not dependent upon its recognition, and
its ontological determinations are not reducible to the
descriptive categories through which they may subsequently
be examined.

The knowing subject is not introduced as a constitutive
principle of the ontology developed here. The subject,
its cognitive capacities, and the conditions under which
it encounters or interprets reality belong to a different
field of philosophical investigation.

This restriction does not imply that knowing subjects
are outside Being. It means only that their existence
and activities must not be used as the starting conditions
for defining the ontological structure of reality.

The distinction between ontology and epistemology is
therefore essential to the architecture of the Philosophy
of Viability. The former investigates Being and its
ontological determinations; the latter concerns the
conditions and structures of knowing.

The epistemological investigation belongs to the
subsequent volume of this philosophical system. The
present book neither presupposes its analytical framework
nor imports its concepts into the construction of the
ontology.

This separation protects the independence of the
ontological inquiry and prevents the conditions under
which reality is known from being confused with the
conditions under which reality exists.

\section*{The Realms of Being}

The second part of this volume investigates the
ontological realms through which Being is differentiated,
structured, and manifested.

The proposed architecture distinguishes four fundamental
realms:

\begin{itemize}
    \item First Determination;
    \item Structural Determination;
    \item Relational Manifestation;
    \item Material Manifestation.
\end{itemize}

Alongside these four fundamental realms, the framework
distinguishes three intermediary domains:

\begin{itemize}
    \item Hub of Decree;
    \item Hub of Signification;
    \item Hub of Realization.
\end{itemize}

The distinction between these two groups is fundamental
to the architecture. The seven domains must not be
interpreted as seven equivalent levels arranged in a
simple linear hierarchy.

The four fundamental realms designate distinct
ontological dimensions within the proposed framework.
The three hubs designate intermediary domains whose
functions and relations must be established through
ontological analysis.

Their intermediary status does not, by itself, establish
that they are merely derivative or ontologically
dependent. Equally, their significance does not justify
treating them as additional fundamental realms without
further argument.

In particular, the Hub of Signification must not be
prematurely reduced to language, subjective meaning,
or psychological representation. Its ontological role
must be examined independently of such reductions.

The relations among the fundamental realms and the
intermediary domains constitute an important part of
the investigation. Their functions, structural
dependencies, and possible modes of interaction cannot
be adequately understood by treating the seven domains
as a simple sequence of stages.

The status and precise position of a possible
Transcendent Ground or Throne-Domain are likewise
not presupposed as settled components of the architecture.
Their inclusion, distinction, or relation to the other
domains requires independent philosophical justification.

\section*{The Question of Light}

The third part addresses Light as an ontological problem.

Light is not introduced merely as a metaphor for knowledge,
awareness, intelligibility, or the disclosure of information.
The ontological question concerning Light must be distinguished
from epistemological questions concerning the conditions
under which something becomes accessible to a knowing subject.

Ontological Light belongs to the investigation of Being
and its manifestation. Epistemic Light, by contrast,
belongs to the investigation of knowing and the conditions
under which something becomes intelligible or accessible
within a cognitive relation.

These two dimensions must not be conflated. An account
of illumination in the context of knowledge cannot,
without further argument, establish the ontological
status of Light.

The third part therefore examines Light within the
ontological architecture of the book. Its purpose is
to clarify the role of Light in relation to Being,
determination, and manifestation without importing
epistemological assumptions into the analysis.

The precise relations between Light and the different
realms of Being must be established through argument
rather than inferred from metaphorical or conventional
uses of the term.

\section*{The Dynamics of Being}

The fourth part investigates the dynamics of Being.

Being cannot be adequately understood through an account
that considers only static determination or observable
change. Its dynamics involve a broader set of ontological
questions concerning becoming, transformation, continuity,
discontinuity, emergence, relation, temporality,
generation, dissolution, and the consequences of
ontological transitions.

These questions cannot be reduced to descriptions of
physical motion or empirical variation. Observable
change may constitute one manifestation of a more
fundamental transformation, but the ontological
conditions of change are themselves a distinct subject
of investigation.

The dynamics of Being also require an examination of
the relations between identity and transformation.
An account of identity that excludes the possibility
of change risks reducing identity to static sameness.
Conversely, an account of change that disregards
identity risks losing the basis upon which the
continuity and distinction of an entity can be examined.

The relation between continuity and discontinuity,
therefore, cannot be settled through an exclusive
emphasis on either persistence or transformation.

Temporality is also approached as an ontological
problem rather than being identified with the
subjective experience or cognitive representation
of time. Its relation to becoming, continuity, and
transformation requires examination within the
structure of Being itself.

The fourth part is consequently not an independent
addition to the preceding discussions. It examines
the dynamic dimensions through which the ontological
architecture must be understood as a whole.

\section*{The Meaning of Viability}

Although this volume is devoted to ontology, its
orientation is determined by the broader philosophical
problem of viability.

Viability must not be identified exclusively with
biological survival, persistence over time, practical
success, adaptation, or the achievement of a
predefined objective. Each of these notions concerns
a particular condition or phenomenon, and none
is sufficient by itself to establish the ontological
meaning of viability.

Nor can viability be reduced to a simple criterion
of truth, functionality, stability, or continued
existence.

The concept requires a more fundamental examination
of the conditions under which entities, identities,
structures, and relations are constituted, maintained,
transformed, or reconfigured within Being.

This examination must distinguish the existence of
an entity from the conditions of its persistence,
the identity of an entity from its particular
manifestations, and the dynamics of its transformation
from the mere observation of change.

Viability is therefore not introduced as a synonym
for survival or persistence. Its precise ontological
meaning must emerge from the systematic examination
of the structures and relations that constitute
the subject of this book.

The present volume does not claim to settle the
complete definition of viability in advance. Rather,
it establishes the ontological questions and
conceptual distinctions that an adequate definition
must address.

This restriction is deliberate. A philosophical
concept must not be defined by presupposing the
very ontological conditions that its investigation
is intended to clarify.

\section*{Established Commitments and Open Questions}

The construction of a systematic philosophy requires
a clear distinction between different degrees and
forms of conceptual commitment.

Throughout this volume, the reader should distinguish
among the following:

\begin{description}
    \item[Established commitments:]
    Foundational positions adopted as constitutive
    commitments of the present ontological framework.

    \item[Provisional formulations:]
    Conceptual constructions that provide a working
    account of an ontological problem but remain
    subject to refinement and further justification.

    \item[Hypotheses:]
    Explanatory or structural proposals whose validity
    has not yet been established.

    \item[Open questions:]
    Problems identified as philosophically significant
    for which no definitive position is presently adopted.

\end{description}

These distinctions are not merely editorial.
They are necessary for maintaining the argumentative
discipline of a developing philosophical system.

Without them, a philosophical framework risks
presenting conjectures as principles, provisional
constructions as established results, or unresolved
problems as if they had already been answered.

The architecture presented in this volume should
therefore be understood as a structured philosophical
undertaking. Its commitments must remain explicit,
its arguments open to examination, and its unresolved
questions distinguishable from its established
positions.

The presence of an open question does not invalidate
the architecture. What matters is whether the
question has been properly identified, whether its
implications are understood, and whether the
framework avoids concealing uncertainty through
premature definitions.

\section*{The Architecture of the Volume}

This book is organized into four parts:

\begin{enumerate}

    \item \textbf{Part I: Identity and Manifestation}

    Examines the ontological problem of identity,
    its inseparable relation to Being, and its
    distinctions from determination, unity,
    distinction, individuation, numerical identity,
    continuity, and signature. It also investigates
    the relation between identity and manifestation
    without reducing identity to any particular
    manifestation.

    \item \textbf{Part II: The Realms of Being}

    Develops the architecture of the four fundamental
    realms and three intermediary hubs through
    which Being is differentiated, structured,
    and manifested. It examines their respective
    ontological functions and relations.

    \item \textbf{Part III: Light}

    Investigates Light as an ontological problem
    and examines its relation to Being and
    manifestation while preserving the distinction
    between ontological and epistemic dimensions.

    \item \textbf{Part IV: Dynamics of Being}

    Examines becoming, change, relation, emergence,
    continuity, discontinuity, temporality,
    transformation, generation, dissolution,
    and the wider conditions and consequences
    of ontological dynamics.

\end{enumerate}

The sequence is conceptually significant.

The investigation begins with identity and manifestation
because the relation between what an entity is and
the ways in which it is determined and manifested
must first be clarified.

It then proceeds to the differentiation of Being
through its fundamental realms and intermediary
domains. The question of Light is subsequently
examined within this ontological architecture.
Finally, the dynamics of Being are investigated
to clarify the processes and transformations
through which the architecture must be understood.

The four parts are therefore intended to function
as connected stages of a single ontological inquiry,
rather than as independent collections of
philosophical topics.

\section*{A Developing Philosophical System}

The Philosophy of Viability is conceived as a
systematic philosophical project. Its development
requires more than the accumulation of definitions,
propositions, or conceptual distinctions.

A coherent system must establish explicit relations
among its foundational commitments, clarify the
status of its concepts, and examine the implications
of its arguments across the different domains of
its architecture.

For this reason, the development of the project
is supported by structured registers of axioms,
definitions, propositions, theses, hypotheses,
arguments, and philosophical foundations.

These registers serve an organizational and
analytical purpose. They facilitate the traceability
of conceptual commitments, the identification of
dependencies among arguments, and the examination
of consistency across the system.

They must not, however, be confused with the
philosophical justification of the claims they
contain.

A proposition does not become justified merely
because it has been registered. A definition
does not become ontologically adequate merely
because it has been formally stated. Nor does
the organization of concepts guarantee the
coherence of their relations.

The registers are therefore instruments of
systematic examination, not substitutes for
philosophical argument.

The framework remains open to criticism,
clarification, and further development. Its
coherence must be established through argument
and examination rather than assumed through
terminology or formal organization.

\section*{Closing Orientation}

This first volume establishes a point of departure,
not a final account of everything that exists.

Its central task is to formulate and investigate
the ontological questions that must be addressed
before a comprehensive philosophy of viability
can be developed.

Identity, determination, manifestation, the
realms of Being, Light, and the dynamics of Being
are approached as interconnected dimensions
of one philosophical problem.

The ambition is not to eliminate the unknown
through definition, nor to force every unresolved
question into a prematurely closed conceptual
system. It is to construct an ontological framework
capable of distinguishing what has been established
from what remains to be investigated.

Such a framework must preserve the distinction
between fundamental and intermediary domains,
between identity and its determinations, between
Being and its manifestations, and between
ontological commitments and provisional proposals.

The Philosophy of Viability begins, therefore,
with the investigation of Being itself. Its
starting point is neither the knowing subject
nor the conditions under which reality becomes
intelligible to that subject.

It begins with the ontological question of what
it means for entities to exist, to possess
identity, to undergo determination, and to
manifest themselves.

This is the foundation upon which the subsequent
development of the Philosophy of Viability
must be constructed.

\cleardoublepage

% ============================================================
% End of Preface
% ============================================================
